Verteilungsbasierte Datensatzsuche adressiert die Aufgabe, Datens\"{a}tze anhand ihrer statistischen Eigenschaften zu finden, wenn die Rohdaten nicht zug\"{a}nglich sind und lediglich Metadaten in Form von spaltenweisen Histogrammen vorliegen.\footnote{Fachbegriffe und ausgew\"{a}hlte Formulierungen dieser Zusammenfassung wurden mit \emph{DeepL} (\url{https://www.deepl.com/}) \"{u}bersetzt, um Terminologie und Satzbau des englischen Abstracts zu bewahren.} Fainder~\cite{behme_2024} ist ein vorgeschlagener Index f\"{u}r dieses Problem: er clustert Histogramme mit K-Means, richtet sie innerhalb jedes Clusters auf ein gemeinsames Bin-Gitter aus und berechnet kumulative Dichten vor, sodass Perzentilpr\"{a}dikate durch Bin\"{a}rsuche anstelle eines vollst\"{a}ndigen Histogrammdurchlaufs beantwortet werden. Auf der GitTables-Kollektion mit f\"{u}nf Millionen Histogrammen erreicht Fainder eine Beschleunigung von mehr als zwei Gr\"{o}\ss{}enordnungen gegen\"{u}ber der \texttt{profile-scan}-Baseline und reduziert den Speicherbedarf von \"{u}ber 1\,TB auf 3{,}2\,GB. Die Referenzimplementierung ist jedoch in Python geschrieben, und der Python-Overhead dominiert ihre Abfragezeit: der Global Interpreter Lock serialisiert die Cluster-Schleife \"{u}ber Threads hinweg, das Array-of-Structs-Layout von NumPy verschmutzt Cache-Lines mit ungelesenen Histogramm-ID-Feldern, und der Interpreter-Dispatch pro Iteration verst\"{a}rkt beides. Das Ergebnis ist eine Laufzeit von $2{,}5$ Stunden f\"{u}r $10\,000$ Abfragen auf dem vollen Workload mit f\"{u}nf Millionen Histogrammen, unabh\"{a}ngig von der Thread-Anzahl.

Diese Arbeit reimplementiert Fainders Abfragephase in Rust, einer Systemprogrammiersprache mit starkem Fokus auf Performance und Speichersicherheit. Die resultierende Engine wird genutzt, um die Leistungsobergrenzen verteilungsbasierter Suche auf moderner Mehrkern-Hardware zu charakterisieren und die Komposierbarkeit der Optimierungen zu untersuchen, die diese Obergrenzen adressieren. Wir finden, dass verteilungsbasierte Suche durch vier verschiedene Leistungsobergrenzen begrenzt ist: Ein-Thread-Rechenzeit / L1-Latenz, gemeinsame L3-Kapazit\"{a}t, Speicherdruck auf Socket-Ebene und der Cross-Socket-UPI-Interconnect. Jede bindet bei einer anderen Thread-Anzahl und wird von einer eigenen Klasse von Interventionen adressiert. Ein weiteres Ergebnis ist die ``negative Komposierbarkeit'' von Optimierungen: unabh\"{a}ngige Ablationen auf f\"{u}nf verschiedenen Achsen zeigen dasselbe Verhalten, eine feinere Optimierung legt sich \"{u}ber eine gr\"{o}bere, die die bindende Obergrenze bereits reduziert hat, und der Overhead der feineren Optimierung erscheint als Netto-Regression statt als abnehmender Ertrag.

Auf der Peak-Scaling-Zelle ($t=16$) erzielt die Rust-Engine eine $530$--$598\times$-Ende-zu-Ende-Beschleunigung gegen\"{u}ber der Python-Baseline auf pr\"{a}zisions- und recall-erhaltenden Ausf\"{u}hrungen. Dieses Verh\"{a}ltnis beruht \"{u}berwiegend darauf, dass Rust die Konstruktion eines Ergebnis-Dictionaries pro \emph{(Abfrage, Cluster)}-Paar in Python vermeidet; die Suchphase selbst erreicht $22{,}7\times$ auf $324$K Histogrammen und f\"{a}llt auf $1{,}34\times$ bei $5{,}02$ Millionen, w\"{a}hrend sich die Arbeitslast den gemeinsamen Hardware-Obergrenzen n\"{a}hert, die diese Arbeit charakterisiert. Weitere $11$--$13\%$ der Wall-Time lassen sich bei $t \leq 48$ durch Bindung des Prozesses an einen einzelnen NUMA-Knoten einsparen. Der ebenfalls nach Rust portierte Rebinning-Kernel tr\"{a}gt eine separate $\sim\!4{,}7\times$-Beschleunigung in der Ausrichtungsphase der Indexkonstruktion f\"{u}r Datens\"{a}tze unter einer Million Histogrammen bei, die bei f\"{u}nf Millionen auf $\sim\!2{,}2\times$ f\"{a}llt (bedingt durch die Parallelisierungsobergrenze des Kernels von ${\sim}2$ effektiven Kernen auf der gemeinsamen Cluster-Allokation). Die Ergebnisse werden in einer Dispatch-Policy konsolidiert, die pro Workload-Regime die beste Build-Variante ausw\"{a}hlt und auf $23$ von $24$ Zellen innerhalb von $5\%$ validiert ist, einschlie\ss{}lich eines Out-of-Distribution-Dichtepunkts. Diese Zahlen sind eine nachgelagerte Konsequenz der obergrenzenbewussten Neuentwicklung und nicht der prim\"{a}re Beitrag dieser Arbeit.

Der methodische Hauptbeitrag ist die \textbf{Pre-Flight-Obergrenzenidentifikation}: bevor eine Optimierung implementiert wird, sondiert man die Zielkomposition mit Hardware-Performance-Countern und schlie\ss{}t jede Optimierung aus, deren Zielachse dort nicht mehr bindet. Das Muster wird f\"{u}r jede Perzentilpr\"{a}dikat-Abfrage-Engine erwartet, deren Zugriffsmuster eine Kette von Bin\"{a}rsuchen pro Cluster ist.
